\begin{solapartat}
\punts{0,50} El camp magnètic màxim a 10~cm $= 0,1$~m el calculem a partir de la intensitat màxima $10^{5}$~A i l'expressió del mòdul del camp magnètic creat per un fil infinit:
\[ B_{max} = \frac{\mu_0 I_{max}}{2\pi r} = \frac{4\pi\cdot 10^{-7}\cdot 10^{5}}{2\pi\cdot 0,1} = 0,2\ \mathrm{T} \]
Dibuix:

\punts{0,10} La transferència de càrrega negativa del núvol cap a terra.

\punts{0,20} La direcció de la intensitat de corrent cap amunt.

\punts{0,20} Les línies de camp són línies concèntriques al voltant d'un fil infinit.

\punts{0,25} El sentit de les línies de camp magnètic al voltant del parallamps ens el dona la regla de la mà dreta o equivalent.

\figura[0.2]{sol_fig1.png}
\end{solapartat}

\begin{solapartat}
\punts{0,75} Gràfica. Calculem alguns punts del camp per representar-lo.

Per exemple: $B(0,1\ \mathrm{m}) = 0,2\ \mathrm{T}$, $B(0,2\ \mathrm{m}) = 0,1\ \mathrm{T}$, $B(0,3\ \mathrm{m}) = 0,0667\ \mathrm{T}$, $B(0,4\ \mathrm{m}) = 0,05\ \mathrm{T}$, $B(0,5\ \mathrm{m}) = 0,04\ \mathrm{T}$

No incloure títol a l'eix resta 0,10 punts.\\
No incloure unitats a l'eix resta 0,2 punts.

\figura[0.4]{sol_fig2.png}

\punts{0,25} La força magnètica sobre l'electró és: $\vec{F} = q\vec{v}\times\vec{B}$. El camp magnètic i la velocitat són perpendiculars i, per tant, el mòdul de la força magnètica serà: $|\vec{F}| = qvB = 1,602\cdot 10^{-19}\cdot 10^{3}\cdot 0,2 = 3,2\cdot 10^{-17}\ N$

\punts{0,25} La regla de la mà dreta (o equivalent) i el signe de la càrrega ens indica que la força serà en la direcció i el sentit cap al parallamps.

\figura[0.2]{sol_fig3.png}
\end{solapartat}
